\documentclass[11pt]{article}
% Shared preamble for the rigor documents (Lamport-style structured proofs).  \input this file after \documentclass.
\usepackage[T1]{fontenc}\usepackage{lmodern}\usepackage{microtype}
\usepackage[margin=1in]{geometry}
\usepackage{amsmath,amssymb,amsthm,mathtools}
\usepackage{xcolor}\usepackage{enumitem}\usepackage{booktabs}\usepackage{graphicx}\usepackage{hyperref}
\usepackage[most]{tcolorbox}
\definecolor{navy}{HTML}{17365D}\definecolor{teal}{HTML}{117A8B}\definecolor{warn}{HTML}{A33A2B}\definecolor{okgreen}{HTML}{2E7D4F}
\hypersetup{colorlinks=true,linkcolor=navy,citecolor=teal,urlcolor=teal}
\theoremstyle{plain}
\newtheorem{theorem}{Theorem}[section]\newtheorem{lemma}[theorem]{Lemma}\newtheorem{proposition}[theorem]{Proposition}\newtheorem{corollary}[theorem]{Corollary}
\theoremstyle{definition}\newtheorem{definition}[theorem]{Definition}\newtheorem{remark}[theorem]{Remark}\newtheorem{claim}[theorem]{Claim}\newtheorem{issue}[theorem]{Issue}
% ---------- Lamport structured proofs ----------
% \lstep{level}{number}{assertion}   -> indented "<level>number. assertion"
% \lproof                             -> "PROOF:" line (start of the sub-proof of the preceding step)
% \lby{justification}                 -> "BY ..." leaf justification for the preceding step
% \lqed{level}{number}                -> QED step of a sub-proof
% keywords: \lassume \lprove \lsuffices \lcase \ldefine \lpick \lnew
\newlength{\lindent}\setlength{\lindent}{1.7em}
\newcommand{\lstep}[3]{\par\smallskip\noindent\hspace*{\dimexpr\numexpr#1-1\relax\lindent\relax}{\color{navy}$\langle#1\rangle#2$.}\ #3}
\newcommand{\lproof}{\par\noindent\hspace*{\lindent}{\scshape Proof:}}
\newcommand{\lby}[1]{\par\noindent\hspace*{\lindent}{\scshape By}\ #1.}
\newcommand{\lqed}[2]{\par\smallskip\noindent\hspace*{\dimexpr\numexpr#1-1\relax\lindent\relax}{\color{navy}$\langle#1\rangle#2$.}\ {\scshape Q.E.D.}}
\newcommand{\lassume}{{\scshape Assume}\ }\newcommand{\lprove}{{\scshape Prove}\ }\newcommand{\lsuffices}{{\scshape Suffices}\ }
\newcommand{\lcase}{{\scshape Case}\ }\newcommand{\ldefine}{{\scshape Define}\ }\newcommand{\lpick}{{\scshape Pick}\ }\newcommand{\lnew}{{\scshape New}\ }
% ---------- verdict boxes ----------
\newtcolorbox{verdictok}{colback=okgreen!8,colframe=okgreen,title={\bfseries Verdict: verified},fonttitle=\small}
\newtcolorbox{verdictgap}{colback=orange!10,colframe=orange!80!black,title={\bfseries Verdict: gap / caveat},fonttitle=\small}
\newtcolorbox{verdictbreak}{colback=warn!10,colframe=warn,title={\bfseries Verdict: proof-breaking issue},fonttitle=\small}
\newtcolorbox{citedbox}[1]{colback=navy!5,colframe=navy!60,title={\bfseries Cited result: #1},fonttitle=\small,breakable}
\setlength{\parindent}{0pt}\setlength{\parskip}{4pt}\allowdisplaybreaks
\newcommand{\cA}{\mathcal A}\newcommand{\cF}{\mathcal F}\newcommand{\cM}{\mathfrak M}\newcommand{\cN}{\mathfrak N}

% extra calligraphic shortcuts (\providecommand: no clash if a document defines its own)
\providecommand{\cB}{\mathcal B}\providecommand{\cC}{\mathcal C}\providecommand{\cD}{\mathcal D}\providecommand{\cE}{\mathcal E}\providecommand{\cG}{\mathcal G}\providecommand{\cH}{\mathcal H}\providecommand{\cI}{\mathcal I}\providecommand{\cJ}{\mathcal J}\providecommand{\cK}{\mathcal K}\providecommand{\cL}{\mathcal L}\providecommand{\cO}{\mathcal O}\providecommand{\cP}{\mathcal P}\providecommand{\cQ}{\mathcal Q}\providecommand{\cR}{\mathcal R}\providecommand{\cS}{\mathcal S}\providecommand{\cT}{\mathcal T}\providecommand{\cU}{\mathcal U}\providecommand{\cV}{\mathcal V}\providecommand{\cW}{\mathcal W}\providecommand{\cX}{\mathcal X}\providecommand{\cY}{\mathcal Y}\providecommand{\cZ}{\mathcal Z}
\newcommand{\Fid}{\operatorname{F}}\newcommand{\Tr}{\operatorname{Tr}}\newcommand{\eps}{\varepsilon}\newcommand{\dd}{\,\mathrm d}

\newcommand{\Wt}{\widetilde W}\newcommand{\kt}{\widetilde k}
\newcommand{\Stwo}{\mathfrak S_2}\newcommand{\Sone}{\mathfrak S_1}
\newcommand{\Pp}{P_+}\newcommand{\Pm}{P_-}\newcommand{\Om}{\Omega}\newcommand{\one}{\mathbf 1}
\newcommand{\norm}[1]{\lVert#1\rVert}\newcommand{\abs}[1]{\lvert#1\rvert}
\newcommand{\ip}[2]{\langle#1,#2\rangle}
\DeclareMathOperator{\dist}{dist}
\title{Referee report on \texttt{rigor/rate\_of\_remainder.tex}\\
\large (track RR: the rate of the remainder; the group bound, $c_3$, and the lower rate)}
\author{Referee RRr}\date{9 September 2026}
\begin{document}\maketitle

\noindent\textbf{Documents refereed.}  \texttt{rigor/rate\_of\_remainder.tex} (agent RR, 15\,pp.),
against \texttt{rigor/uhlmann\_upper\_bound.tex} (PA), \texttt{rigor/quadratic\_limit.tex} (P2),
\texttt{fidelity\_recovered\_quasifree.tex} (Note~5), \texttt{fidelity\_tables.tex},
\texttt{rigor/check\_note5\_theory.tex}, \texttt{numerics/p1\_bogoliubov.py/.out}.  All numerical
re-checks below were run independently with
\texttt{scratchpad/rrr\_check.py} (numpy/mpmath).

\section*{Executive summary}
\begin{center}\small
\begin{tabular}{p{0.36\textwidth}p{0.30\textwidth}l}\toprule
ingredient & my finding & verdict\\\midrule
(a1) flow extension is a one-parameter group; generator $=$ PA's $D$ &
 correct; $h_s\circ h_t=h_{s+t}$ verified; ``for all $s$'' should read ``for all $s\ge0$'' &
 verified (MINOR)\\
(a2) $\norm{V_s}_2^2\le s^2\norm{\Pm D\Pp}_2^2$ exactly &
 correct; the autocorrelation apparatus is unnecessary (Bochner triangle inequality suffices) and
 hides that the result holds for \emph{every} one-parameter unitary group & verified\\
(a3) $\Phi\le-\frac12\log(1-2f_2\zeta^2)$, i.e.\ $f_2^\star=f_2$ &
 RR's own class of flow extensions leaves this open, as RR says; but the gap is
 \textbf{closable}: enlarging the class to $U_s=e^{s(D+ih')}$, $h'\in\cA'$, keeps the group property,
 keeps the locality, and makes PA Prop.~7.6 give $f_2^\star=f_2$ \emph{exactly} & gap, with a repair\\
(a4) evenness of $\det(\one-V_s^*V_s)$ & a theorem (CS decomposition), not numerics; and it is not
 needed for the conclusion & verified\\
(b1) third-order formula $\Phi_3$ & re-derived and verified numerically on random $3$- and
 $4$-mode states, four independent draws & verified\\
(b2) $c_3=-0.99(1)$, $c_4=+0.9(1)$ & reproduced; and I find the \emph{published}
 Table~\ref{tab:tab3} (Note~5 Tab.~3, col.~A) is matched to $2\cdot10^{-6}$ relative for
 $\zeta\le\frac1{30}$ --- stronger than RR claims & verified (numerical)\\
(c) $C(\Lambda)$ bounded $\Rightarrow$ linear lower rate & the boundedness is measured for the
 Taylor coefficient only; Problem~7.2 needs a uniform third derivative.  RR says so.  &
 gap (correctly declared)\\
(d) refutation of the $\log^{-2}$ remainder of N5 Cor.~4.2 & logically complete
 \emph{as a refutation of that remainder as stated}, unconditionally; see \S\ref{sec:logic} for the
 precise scope & verified\\
\bottomrule\end{tabular}\end{center}

\textbf{No proof-breaking issue was found in \texttt{rate\_of\_remainder.tex}.}  The
one \texttt{verdictbreak} the document itself issues (against N5 Cor.~4.2 and against P2
Problem~7.2's guess $C(\Lambda)\sim e^\Lambda$) is, in my judgement, justified in the first case
and is a numerical --- not proved --- statement in the second; \S\ref{sec:c} says exactly what is
and is not established.  The single substantive criticism is that RR stops one step short: the
conditional clause ``modulo $f_2^\star=f_2$'' that decorates the main theorem and
Corollary~2.7 (\texttt{cor:c3sign}) can be removed, by the argument of \S\ref{sec:repair}.

\tableofcontents

\section{(a1) Is the map on $I$ the flow of $\chi=x^2$ for all $s$ at once?}\label{sec:a1}

\textbf{Yes, for $s\ge0$; and this is PA's construction, not a new one.}  Three separate claims must
be distinguished, and RR conflates them slightly.

\paragraph{(i) The M\"obius family is a group.}  With $h_s(x)=Lx/(L+sx)$,
\[
 h_s(h_t(x))=\frac{L\cdot\frac{Lx}{L+tx}}{L+s\frac{Lx}{L+tx}}
 =\frac{L^2x}{L(L+tx)+sLx}=\frac{Lx}{L+(s+t)x}=h_{s+t}(x),
\]
an identity of rational functions valid whenever no denominator vanishes.  Verified numerically at
$(s,t,x)=(0.3,0.5,0.7),(1.3,-0.4,0.9),(0.05,0.05,1)$ to $15$ digits.  Equivalently
$h_s=\sigma^{-1}\circ\tau_s\circ\sigma$ with $\sigma(x)=L/x$ and $\tau_s(y)=y+s$; and
$\partial_sh_s(x)=-h_s(x)^2/L$, so with the normal form $L=1$ the family is the flow of $-x^2\partial_x$.

\paragraph{(ii) The orbit stays in $[0,1]$ only for $s\ge0$.}  RR writes (Def.~2.1 ff.)
``\emph{an admissible field produces the required recovery map on $I$ for the whole one-parameter
family at once}'' and, in Theorem~2.3, ``for every $s\in\mathbb R$''.  For $s<0$ the flow of
$\chi\ge0$ moves points to the \emph{right}, so a point of $[0,1]$ leaves $[0,1]$ and enters the
region where $\chi\ne x^2$; then $\kt_s\ne h_s$ on $[0,1]$ (indeed $h_{-1/x}(x)=\infty$).  PA states
this correctly: Lemma~2.2\,(E2) is asserted ``for every $s\ge0$''.

\begin{verdictgap}
Cosmetic but worth fixing.  Theorem~2.3 of \texttt{rate\_of\_remainder.tex} is a statement about the
unitary group $\Wt_s=e^{sD}$ alone and is true for all $s\in\mathbb R$; the sentence after
Definition~2.1 (``\emph{simultaneously for all $s$}'') is an assertion about the geometry and is true
only for $s\ge0$.  Since only $\zeta=s\ge0$ is ever used, nothing downstream changes; insert
``$s\ge0$'' in the sentence after (F2).
\end{verdictgap}

\paragraph{(iii) The extension, the group property, and the kink.}  $\chi$ is Lipschitz with
$\chi'$ Lipschitz, so Picard--Lindel\"of gives a unique global flow and hence
$\kt_s\circ\kt_t=\kt_{s+t}$ on all of $\mathbb R$; the jump of $\chi''$ at $x=0$ (of size $2$, PA
Lemma~2.2) is irrelevant --- uniqueness of solutions needs only local Lipschitz continuity of
$\chi$, and $\kt_s'$ is given by the exponential formula PA Lemma~2.2\,(E3) with a merely
$L^\infty$ integrand.  Since $W_\phi W_\psi=W_{\phi\circ\psi}$ (PA eq.~(3)), $s\mapsto\Wt_s$ is a
one-parameter unitary group; PA already proves this and computes the generator, PA Lemma~5.1:
\[
 \Wt_s=e^{sD},\qquad D=\chi\partial_x+\tfrac12\chi',\qquad D^*=-D,\quad H^1_c\ \text{a core}.
\]
This is \emph{verbatim} RR's $D$ in (2.2).  So the answer to ``is the generator the same as PA's''
is yes, identically; and the group property is not new information --- PA Lemma~5.1
$\langle1\rangle1$ states it.  What is new in RR is only the observation that the group property
alone already collapses PA's error term, see \S\ref{sec:a2}.

\begin{verdictok}
(a1) is correct.  The one-parameter group $\Wt_s=e^{sD}$ with $D=\chi\partial_x+\frac12\chi'$ is
PA's; the flow reproduces $k_s$ on $I$ for every $s\ge0$; the $C^{1,1}$ kink at $x=0$ (and any kink
placed at $x=1$ by the extension) does not disturb the group law.  The paper should say plainly that
Lemma~2.2 and Theorem~2.3 are statements \emph{about PA's own family with $h'=0$}, not about a new
construction.
\end{verdictok}

\begin{verdictgap}
Hypotheses inherited from PA that RR's class violates.  RR's admissible fields are \emph{not}
required to be compactly supported --- and must not be, since the near-optimal ones are
asymptotic translations, $\chi\to\text{const}\ne0$ as $x\to+\infty$
(\texttt{p1\_bogoliubov.out}~(2): ``\emph{the optimal extension moves the point at infinity}\ldots
restricting to $u\to0$ gives a factor $8.0$ too large'').  But PA Lemma~2.4 and PA Prop.~2.5 --- the
sources for the kernel formula and for $G_\phi\in\Stwo$ --- both assume $\phi=\mathrm{id}$ outside a
compact set.  RR's (F3) ($\kappa^{(1)}\in L^2$) is the right substitute and does give
$\kappa_{\kt_s}\in L^2$ through Lemma~2.2, but the paper should record that PA Lemma~2.4/Prop.~2.5
must be re-run for diffeomorphisms that are asymptotically translations (the kernel computation is
unaffected; only the $L^2$ bookkeeping is).  The same remark applies to Lemma~2.8
$\langle1\rangle3$ $\langle2\rangle2$, which argues ``for $y$ outside $\operatorname{supp}\chi$''
and therefore covers only compactly supported fields.
\end{verdictgap}

\section{(a2) The exact identity $\norm{V_s}_2^2\le s^2\norm{\Pm D\Pp}_2^2$}\label{sec:a2}

I re-derived Theorem~2.3 of \texttt{rate\_of\_remainder.tex} from scratch.  It is correct.  It is
also considerably more elementary --- and considerably more general --- than the proof given.

\begin{theorem}[what is actually true]\label{thm:mine}
Let $(U_s)_{s\in\mathbb R}$ be \emph{any} strongly continuous one-parameter unitary group on
$H=L^2(\mathbb R)$ with skew-adjoint generator $K$ such that $[\Pm,K]\in\Stwo$ and
$s\mapsto U_s^*\Pp U_s$ is $\Stwo$-differentiable.  Put $V_s:=\Pm U_s\Pp$,
$G_s:=\Pp-U_s^*\Pp U_s$, $G^{(1)}:=[\Pm,K]$.  Then for all $s\in\mathbb R$
\begin{equation}\label{eq:mine}
 \norm{V_s}_2^2=\tfrac12\norm{G_s}_2^2\le\tfrac12s^2\norm{G^{(1)}}_2^2=s^2\norm{\Pm K\Pp}_2^2 ,
\end{equation}
and $\norm{V_{-s}}_2=\norm{V_s}_2$, indeed $\det(\one-V_{-s}^*V_{-s})=\det(\one-V_s^*V_s)$.
\end{theorem}
\lstep{1}{1}{For any unitary $U$ and $V=\Pm U\Pp$, $B=\Pp U\Pm$:
 $\norm{\Pp-U^*\Pp U}_2^2=\norm V_2^2+\norm B_2^2$.}
\lby{$P:=\Pp$, $Q:=U^*\Pp U$ are orthogonal projections, and
 $\Tr[(P-Q)^2]=\Tr[P(\one-Q)]+\Tr[Q(\one-P)]=\norm{(\one-Q)P}_2^2+\norm{(\one-P)Q}_2^2$;
 $(\one-Q)P=U^*\Pm U\Pp$ has the same $\Stwo$-norm as $\Pm U\Pp=V$, and
 $(\one-P)Q=\Pm U^*\Pp U$ has the same $\Stwo$-norm as $\Pp U\Pm=B$}
\lstep{1}{2}{$V$ and $B$ have the same singular values, so $\norm B_2=\norm V_2$ and
 $\det(\one-B^*B)=\det(\one-V^*V)$.}
\lby{unitarity of $U$ gives $V^*V=\one_{\Pp H}-A^*A$ and $BB^*=\one_{\Pp H}-AA^*$ with
 $A=\Pp U\Pp$; $A^*A$ and $AA^*$ have the same spectrum away from $0$, hence so do $V^*V$ and
 $BB^*$; equivalently, the CS decomposition of $U$ makes both off-diagonal blocks equal to the same
 $\sin\Theta$}
\lstep{1}{3}{$G_s=\int_0^sU_\sigma^*G^{(1)}U_\sigma\dd\sigma$ as a Bochner integral in $\Stwo$.}
\lby{$\frac{\dd}{\dd\sigma}U_\sigma^*\Pp U_\sigma=U_\sigma^*[\Pp,K]U_\sigma=-U_\sigma^*G^{(1)}U_\sigma$
 (product rule on the core; $[\Pp,K]=-[\Pm,K]$ because $\Pp+\Pm=\one$), and
 $\sigma\mapsto U_\sigma^*G^{(1)}U_\sigma$ is $\Stwo$-continuous with constant norm
 $\norm{G^{(1)}}_2$ (PA Lemma~5.2 $\langle1\rangle4$, whose proof uses only strong continuity of
 the group and $\Stwo$-density of finite-rank operators)}
\lstep{1}{4}{$\norm{G_s}_2\le\abs s\,\norm{G^{(1)}}_2$.}
\lby{$\langle1\rangle3$ and the triangle inequality for the Bochner integral,
 $\norm{\int_0^sF_\sigma\dd\sigma}_2\le\int_0^s\norm{F_\sigma}_2\dd\sigma$, together with
 $\norm{U_\sigma^*G^{(1)}U_\sigma}_2=\norm{G^{(1)}}_2$}
\lstep{1}{5}{$\norm{G^{(1)}}_2^2=2\norm{\Pm K\Pp}_2^2$.}
\lby{PA Lemma~5.2 $\langle1\rangle5$: $[\Pm,K]=\Pm K\Pp-\Pp K\Pm$, the two blocks are
 $\Stwo$-orthogonal, and $\Pp K\Pm=-(\Pm K\Pp)^*$ since $K^*=-K$}
\lqed{1}{6}\lby{$\langle1\rangle1$--$\langle1\rangle5$}

\paragraph{Comparison with RR's proof.}  RR routes $\langle1\rangle4$ through the autocorrelation
function $R(u)=\ip{F_0}{F_u}$ (its Lemma~2.2 and Theorem~2.3 $\langle1\rangle2$).  Every step of
that route is correct: (i) $F_\sigma$ is the kernel of $-2\pi i\,\Wt_\sigma^*G^{(1)}\Wt_\sigma$ ---
I checked the conjugation rule, $W_\phi^*AW_\phi$ has kernel
$\sqrt{\phi'(x)\phi'(y)}\,a(\phi(x),\phi(y))$, and the Jacobian identity that makes
$\norm{F_\sigma}_{L^2}$ constant; (ii) stationarity,
$\ip{F_\sigma}{F_\tau}=(2\pi)^2\Tr[G^{(1)*}\Wt_uG^{(1)}\Wt_u^*]$ with $u=\sigma-\tau$, which is
exactly where the group law $\Wt_\sigma\Wt_\tau^*=\Wt_{\sigma-\tau}$ enters and is the only place it
is needed; (iii) $R$ real and even because $\kappa^{(1)}$ is real and symmetric; (iv)
$\abs{R(u)}\le\norm{F_0}\norm{F_u}=R(0)$ by Cauchy--Schwarz.  But (iv) \emph{is} the Bochner
triangle inequality in disguise: $\int_0^s\!\int_0^sR(\sigma-\tau)=\norm{\int_0^sF_\sigma}^2\le
(\int_0^s\norm{F_\sigma})^2=s^2R(0)$ is the same two lines.  The autocorrelation picture earns its
keep only in Corollary~2.9, where $R(0)-R(u)=\frac12\norm{F_u-F_0}^2$ gives the \emph{quartic}
coefficient; there it is genuinely the right tool.

\begin{verdictok}
Theorem~2.3 and Lemma~2.2 of \texttt{rate\_of\_remainder.tex} are correct as stated and as proved.
The chain of normalisations checks out: PA eq.~(22) gives
$\norm{\Pm D\Pp}_2^2=\frac12\norm{G^{(1)}}_2^2=\frac1{8\pi^2}\iint(\kappa^{(1)})^2$, so RR's
$f_2^\chi=\frac12\norm{\Pm D\Pp}_2^2=\frac1{16\pi^2}\iint(\kappa^{(1)})^2$ agrees with PA
Cor.~5.5, in which $\limsup\Phi/\zeta^2\le\frac12\norm{\Pm(D+ih')\Pp}_2^2$.  The identity
$E(s)=\frac1{8\pi^2}\iint\kappa_{\kt_s}^2$ used in Theorem~2.3 $\langle1\rangle1$ is exact, not
asymptotic, for the reason given in $\langle1\rangle1$--$\langle1\rangle2$ above (RR cites
\texttt{p1\_bogoliubov.out}~(1) for it, which is legitimate but is a numerics file; the two-line
proof should be in the paper).
\end{verdictok}

\section{(a3) The global bound, $f_2^\star$, and how to remove the hypothesis}\label{sec:repair}

\subsection{What is proved and what is numerical, exactly}

\begin{itemize}\itemsep1pt
\item \textbf{Proved, unconditionally, for each fixed admissible field $\chi$:}
 $\Phi(\zeta)\le-\frac12\log(1-2f_2^\chi\zeta^2)$ for $2f_2^\chi\zeta^2<1$
 (Thm.~2.5).  I checked each step.  $\langle1\rangle1$ needs PA Prop.~3.4 (Uhlmann) and PA Thm.~4.1
 (Bogoliubov overlap), whose hypotheses are $\norm V<1$ \emph{and} $\norm{\Pp U\Pm}<1$
 (PA Cor.~4.3 and the \texttt{citedbox} preceding it); RR quotes only the first, but by
 Theorem~\ref{thm:mine} $\langle1\rangle2$ the two norms are equal, so nothing is missing --- worth a
 half-sentence.  $\langle1\rangle2$, $\prod(1-a_j)\ge1-\sum a_j$, is correct.
\item \textbf{Proved:} $f_2^\star\ge f_2$ (Prop.~2.6 $\langle1\rangle1$).  The argument is clean and
 does not use any subspace theory: let $\zeta\downarrow0$ in the bound and use PA Thm.~8.2
 ($\lim\Phi/\zeta^2=f_2$).
\item \textbf{Numerical:} $f_2^\star\le0.00844939$ ($=1.00071f_2$) at basis size $M=200$, with a
 $1/M$ Richardson limit $0.00844649$ ($=1.00036f_2$).  Source \texttt{p1\_bogoliubov.out}~(4).  I
 checked the normalisation of the factor $4.5$: p1 works at $a=1$, $L=2$, so $\zeta=s/3$ and
 $\frac12E/\zeta^2=\frac12\cdot9\,Q=4.5\,Q$ with $Q[g]=\frac1{48\pi^2}\int_0^\infty k^3\abs{\hat
 g}^2$; this is consistent with RR's own $Q=2f_2^\chi$ in \S2.4, where the normal form has
 $\zeta=s$.  Consistent.
\item \textbf{Therefore conditional:} $\Phi\le-\frac12\log(1-2f_2\zeta^2)$, and with it $c_3\le0$
 (Cor.~2.7 $\langle1\rangle1$) and the clean refutation of a $+\zeta^2/\log^2(1/\zeta)$ remainder.
 The unconditional fallback (Cor.~2.7 $\langle1\rangle2$) excludes N5's term only for
 $\zeta\ge6\cdot10^{-17}$ (or $2\cdot10^{-23}$ with Richardson).  I recomputed both thresholds:
 $1/(120\cdot7.1\cdot10^{-4}f_2)=1390$, $e^{-\sqrt{1390}}=6.3\cdot10^{-17}$, and
 $1/(120\cdot3.6\cdot10^{-4}f_2)=2742$, $e^{-\sqrt{2742}}=1.8\cdot10^{-23}$.  Correct.
\end{itemize}

RR's diagnosis of the gap (the \texttt{verdictgap} after Prop.~2.6) is accurate as far as it goes:
PA optimises over \emph{bounded} $h'$ localised in $I'$, RR over vector fields, i.e.\ over the
\emph{unbounded} skew-adjoint generators $D_\eta=\eta\partial_x+\frac12\eta'$ with
$\operatorname{supp}\eta\subseteq I'$, and there is no reason for the two infima to coincide.  But
RR frames this as needing a density theorem in $H_{\cM'}$.  It does not.  The correct move is to
notice that RR's own theorem never used the geometric nature of the family.

\subsection{The repair: use the group generated by $K=D+ih'$}

\begin{proposition}[$f_2^\star=f_2$, unconditionally]\label{prop:repair}
For $h'\in\cA'$ (PA Def.~5.3: $h'=h'^*$ bounded, $h'=E_{I'}h'E_{I'}$, $\Pm h'\Pp\in\Stwo$) put
$K:=D+ih'$ and $U_s:=e^{sK}$.  Then for every $s\ge0$ the vector $\Gamma(U_s)^*\Om$ is a
representative of $\widetilde\omega_s$ on $\cM=\cF(I)$, and consequently
\[
 \Phi(\zeta)\ \le\ -\tfrac12\log\bigl(1-2f_2^K\zeta^2\bigr),\qquad
 f_2^K:=\tfrac12\norm{\Pm K\Pp}_2^2 ,\qquad
 \inf_{h'\in\cA'}f_2^K=\frac{g_B}8=f_2 .
\]
Hence $\Phi(\zeta)\le-\frac12\log(1-2f_2\zeta^2)$ for all $\zeta<(2f_2)^{-1/2}$, with no
hypothesis.
\end{proposition}
\lstep{1}{1}{$U_s=\Wt_s\,c_s$ with $c_s=\mathcal T\exp\bigl(i\int_0^sh'_\sigma\dd\sigma\bigr)$,
 $h'_\sigma:=\Wt_\sigma^*h'\Wt_\sigma$.}
\lby{$c_s:=\Wt_s^{-1}U_s$ satisfies $\partial_sc_s=\Wt_s^{-1}(ih')U_s=(ih'_s)c_s$, $c_0=\one$; the
 time-ordered exponential of a norm-continuous bounded family solves it uniquely}
\lstep{1}{2}{$h'_\sigma=E_{J_\sigma}h'_\sigma E_{J_\sigma}$ with
 $J_\sigma=(\kt_{-\sigma}(1),\infty)\subseteq I'$ for every $\sigma\ge0$.}
\lby{$W_\phi^*(m_f)W_\phi=m_{f\circ\phi}$ for the half-density representation, so
 $\Wt_\sigma^*E_{I'}\Wt_\sigma=E_{\kt_\sigma^{-1}(I')}$ and
 $\kt_\sigma^{-1}\bigl((1,\infty)\bigr)=(\kt_{-\sigma}(1),\infty)$; and $\kt_{-\sigma}(1)\ge1$ for
 $\sigma\ge0$ because $\chi\ge0$}
\lstep{1}{3}{$\Pm h'_\sigma\Pp\in\Stwo$ for every $\sigma$, so each $e^{i\delta h'_\sigma}$ is
 implementable and $\Gamma(c_s)\in\cF(I')$ (twisted, as in PA Lemma~3.2).}
\lby{insert $\one=\Pp+\Pm$ twice in $\Pm\Wt_\sigma^*h'\Wt_\sigma\Pp$: each of the four terms carries
 a factor $\Pm\Wt_\sigma^*\Pp$, $\Pm h'\Pp$, or $\Pm\Wt_\sigma\Pp$, all in $\Stwo$ (PA Prop.~2.5,
 PA Def.~5.3); $\Gamma$ of a norm-limit of products of unitaries localised in $J_\sigma\subseteq I'$
 lies in the von Neumann algebra generated by them}
\lstep{1}{4}{$\Gamma(U_s)^*\Om$ represents $\widetilde\omega_s$ on $\cM$.}
\lby{$\langle1\rangle1$--$\langle1\rangle3$ and PA Prop.~3.1 verbatim, with $\widehat W'$ replaced by
 $c_s$: the commutant factor drops out of $\ip{\Gamma(U_s)^*\Om}{m\,\Gamma(U_s)^*\Om}$}
\lstep{1}{5}{The bound.}
\lby{Theorem~\ref{thm:mine} applied to the group $U_s$ (its hypotheses hold: $[\Pm,K]\in\Stwo$ by
 $\langle1\rangle3$ and PA Lemma~5.2), then PA Prop.~3.4, PA Thm.~4.1 and
 $\prod(1-a_j)\ge1-\sum a_j$ as in Thm.~2.5 $\langle1\rangle1$--$\langle1\rangle3$}
\lqed{1}{6}\lby{$\langle1\rangle5$, PA Prop.~7.6
 ($\inf_{h'\in\cA'}\norm{\Pm(D+ih')\Pp}_2^2=g_B/4$), and continuity of
 $t\mapsto-\frac12\log(1-2t\zeta^2)$ so that the infimum passes to the limit}

\begin{verdictgap}
\textbf{This is the one substantive request to the author.}  Proposition~\ref{prop:repair} says that
the restriction to \emph{flow} extensions --- the source of the whole $f_2^\star$ difficulty --- is
gratuitous.  The only property Theorem~2.3 uses is that $s\mapsto U_s$ is a one-parameter unitary
group whose generator has a Hilbert--Schmidt off-diagonal block; and the family
$U_s=e^{s(D+ih')}$, $h'\in\cA'$, is such a group \emph{and} an admissible purification family,
because the Dyson cocycle $c_s=\Wt_s^{-1}U_s$ stays localised in $I'$ (the flow shrinks $I'$ for
$s\ge0$).  With PA's own Prop.~7.6 the infimum is then $g_B/8=f_2$ \emph{exactly}, and the main
theorem, Corollary~2.7 and the corrected N5 Cor.~4.2 all become unconditional.

Two points still need care before this is written up as a theorem.  (i) The
$\Stwo$-differentiability of $\sigma\mapsto U_\sigma^*\Pp U_\sigma$ that Theorem~\ref{thm:mine}
$\langle1\rangle3$ assumes: for the pure flow it is PA Lemma~5.2 (an explicit kernel computation);
for $K=D+ih'$ one has $\frac{\dd}{\dd\sigma}U_\sigma^*\Pp U_\sigma=-U_\sigma^*[\Pm,K]U_\sigma$
weakly on the core $H^1_c\cap D(h')$, and one must upgrade weak to Bochner --- routine, since the
integrand is $\Stwo$-continuous with constant norm, but it is a step.  (ii) PA Prop.~3.1 and
Lemma~3.2 are stated for $U=\Wt_s\widehat W'$ with $\widehat W'$ a \emph{single} localised unitary;
the time-ordered $c_s$ is a norm limit of finite products of such, so the extension is by continuity
of $\Gamma$ on the relevant Shale group, which should be spelled out.  Neither point looks
problematic; both are cheaper than a density theorem in $H_{\cM'}$.
\end{verdictgap}

\begin{verdictok}
Failing the repair, the document's own accounting is honest and I endorse it: Theorem~2.5 is proved
for each admissible $\chi$; $f_2^\star\ge f_2$ is proved; $f_2^\star\le1.00036\,f_2$ is numerical
(a genuine minimising sequence of admissible fields, converging like $1/M$ in a complete basis, in
a linear least-squares problem --- as strong as numerical evidence gets); the conditional
statements are labelled as such throughout.
\end{verdictok}

\section{(a4) Evenness of $\det(\one-V_s^*V_s)$}\label{sec:even}

\begin{verdictok}
It is a \emph{theorem}, not a numerical observation, and RR proves it correctly: it is
Theorem~\ref{thm:mine} $\langle1\rangle2$ above, i.e.\ the CS decomposition (or, without CS: unitarity
gives $V_s^*V_s=\one-A^*A$ and $B_sB_s^*=\one-AA^*$ with $A=\Pp\Wt_s\Pp$, and $A^*A,AA^*$ are
isospectral off $0$), combined with $V_{-s}=\Pm\Wt_s^*\Pp=(\Pp\Wt_s\Pm)^*$.  The numerical table in
\S2.4 confirms $E(-s)=E(s)$ to $7$--$8$ digits, which is a check of the quadrature, not of the
mathematics.  It should be stated as valid for \emph{any} unitary, since the proof uses nothing else.

One editorial point: evenness is \emph{not} used anywhere in the chain that matters.  The conclusion
``no cubic term in the bound'' is immediate from the closed form $-\frac12\log(1-2f_2^\chi s^2)$,
which is even by inspection.  Evenness of the exact determinant is a nice remark about the
\emph{sharp} group bound $-\frac12\log\det(\one-V_s^*V_s)$; but note that this sharp bound majorises
$\Phi(\zeta)$ only for $\zeta=s\ge0$ (for $s<0$ the map $\kt_s$ is not a recovery map, \S\ref{sec:a1}),
so evenness alone would not license reading off the cubic coefficient of $\Phi$.  The paper's logic
does not do that, but the abstract's phrase ``\emph{with $\det(\one-V_s^*V_s)$ even in $s$
(no cubic term)}'' invites the misreading.
\end{verdictok}

\section{(b) The third-order formula and the value of $c_3$}\label{sec:b}

\subsection{The formula}

I re-derived Theorem~3.1 independently from N5 Thm.~3.2 /
\texttt{check\_note5\_theory.tex} Thm.~2.1,
$\Fid=\det[(1-Q_1)(1-Q_2)]^{1/2}\det[\one+(G_1^{1/2}G_2G_1^{1/2})^{1/2}]$,
$G_j=Q_j(1-Q_j)^{-1}$, and reproduce every line:
\begin{itemize}\itemsep1pt
\item $(1-Q-\eps D)^{-1}=\sum_{n\ge0}\eps^nC^{-2}(DC^{-2})^n$ and $G_\eps=(1-Q_\eps)^{-1}-\one$ give
 $T=G^2+\sum_{j\ge1}\eps^jA_j$, $A_j=G^{1/2}C^{-2}(DC^{-2})^jG^{1/2}$;
\item squaring $T^{1/2}=G+\eps T_1+\eps^2T_2+\eps^3T_3$ gives exactly the Sylvester chain (3.1);
 in the eigenbasis of $Q$, $(T_1)_{ik}=(A_1)_{ik}/(g_i+g_k)$ with
 $(A_1)_{ik}=\sqrt{g_ig_k}D_{ik}/(c_ic_k)$ and $g_i+g_k=N_{ik}/(c_ic_k)$, so
 $(T_1)_{ik}=\sqrt{g_ig_k}D_{ik}/N_{ik}$ and $(M_1)_{ik}=c_i\sqrt{g_ig_k}D_{ik}/N_{ik}$, as printed;
\item $(\one+G)^{-1}=\one-Q=C^2$, so $M=C^2(T^{1/2}-G)$ and
 $-\log\det(\one+T^{1/2})=-\log\det(\one+G)-\Tr M+\frac12\Tr M^2-\frac13\Tr M^3+O(\eps^4)$;
\item $\eps^0$: $-\log\det C^2-\log\det C^{-2}=0$.  $\eps^1$:
 $\Tr M_1=\frac12\sum_iD_{ii}/c_i=\frac12\Tr(C^{-2}D)$ cancels the $\eps^1$ term of
 $-\frac12\log\det(1-Q_\eps)$;
\item $\eps^3$: $\frac16\Tr[(C^{-2}D)^3]-\Tr M_3+\Tr(M_1M_2)-\frac13\Tr M_1^3$, which is (3.3).
\end{itemize}
The bookkeeping of the $\eps^3$ terms of $-\Tr M+\frac12\Tr M^2-\frac13\Tr M^3$ is right:
$-\Tr M_3$, then $\frac12\Tr(M_1M_2+M_2M_1)=\Tr(M_1M_2)$, then $-\frac13\Tr M_1^3$.

\paragraph{Independent numerical test (\texttt{scratchpad/rrr\_check.py}).}  Random
$q_i=1/(1+e^{2g_i})$, $g_i\sim N(0,1)$, random Hermitian $D$ with $\norm D_F=1$, exact $\Fid$ from
the determinant formula; the table lists $(-\log\Fid-\Phi_2\eps^2)/\eps^3$ and the residual divided
by $\eps$ (which must tend to the \emph{quartic} coefficient, i.e.\ stay bounded).
\begin{center}\small
\begin{tabular}{llrrrr}\toprule
$n$ & draw & $\Phi_3$ (formula) & $\eps=10^{-2}$ & $10^{-3}$ & residual$/\eps$ at $10^{-3}$\\\midrule
3 & 0 & $-0.8100206546$ & $-0.784168$ & $-0.807363$ & $+2.658$\\
3 & 1 & $+22.6920764508$ & $+27.574741$ & $+23.091981$ & $+399.9$\\
4 & 0 & $-8.8724102440$ & $-8.075946$ & $-8.784790$ & $+87.62$\\
4 & 1 & $+3.8258778983$ & $+4.211951$ & $+3.861350$ & $+35.47$\\
\bottomrule\end{tabular}\end{center}
In all four draws the difference quotient converges to $\Phi_3$ with residual exactly $O(\eps)$
(the last column is stable between $\eps=10^{-3}$ and $3\cdot10^{-4}$ to a few percent).  Note that
this test validates $\Phi_2$ as well, and much more sharply than a direct comparison: an error in
$\Phi_2$ of relative size $\delta$ would make the third column blow up like $\delta/\eps$.  It does
not.  The second-order part therefore \emph{is} the quantum-Fisher weight
$\frac14\sum\abs{D_{ik}}^2/N_{ik}$ of N5 Prop.~2.3; on the diagonal it is
$\abs{D_{ii}}^2/(8q_ic_i)$, i.e.\ $\frac18\times$ the classical Fisher information of a Bernoulli
variable, the standard Bures normalisation.

\begin{verdictok}
Theorem~3.1 of \texttt{rate\_of\_remainder.tex} is correct, and the derivation is complete
(every step is a genuine step; the only implicit hypothesis, $\eps\norm{C^{-2}D}<1$ and $G>0$, is
stated).  Verified numerically on four independent random $3$- and $4$-mode examples against the
exact determinant formula.  The finite-dimensional restriction is harmless for the application,
where everything is done inside a finite window.
\end{verdictok}

\subsection{The numerical value $c_3=-0.99(1)$: systematics}\label{sec:c3num}

\paragraph{How $\mathcal D_1,\mathcal D_2$ are obtained.}  RR fits
$\widehat D(s)=\sum_{k=1}^4s^kA_k$ exactly through the four stored matrices
$s=0.025,0.05,0.1,0.2$ and sets $\mathcal D_1=-3A_1$, $\mathcal D_2=-9A_2$ (the factors $3,9$ are
the $\zeta=s/3$ conversion of the $a=1,L=2$ configuration; consistent with the rest of the
document).  This is an interpolation, not a least squares, so the systematic error is the
$s^5$ tail of $\widehat D$; RR reports that a cubic fit on $s\le0.1$ moves the results by
$<10^{-5}$, which is the right convergence test and is adequate.  A cleaner alternative --- worth a
sentence in the paper --- would be Richardson-style central differences,
$\mathcal D_1\propto\widehat D(s)-\widehat D(-s)$, which kills the even orders outright.

\paragraph{Box versus window.}  The two systematics are cleanly separated and correctly attributed:
$\Lambda_{\text{box}}=60$ (with $\kappa_{\max}=30$, $N=91$) versus $120$ ($\kappa_{\max}=45$,
$N=273$) move $c_3^{(P)}$ by at most $3\cdot10^{-3}$ relative at fixed $\kappa_c$, whereas the
$\kappa_c$-dependence runs from $-1.098$ to $-1.005$.  So the extrapolation is a
\emph{window} extrapolation, as claimed.  I reproduce RR's own extrapolation of the column
$C(\kappa_c)$: a two-point $1/\kappa_c^2$ Richardson on $(\kappa_c=20.72,23.03)$ gives
$C_\infty=0.008381$, on $(18.42,23.03)$ gives $0.008376$; hence
$c_3=-C_\infty/f_2=-0.9926,\,-0.9920$.  The quoted $C_\infty=0.00838(3)$, $c_3=-0.99\pm0.01$ is
therefore \emph{conservative}, not optimistic: the internal spread of the extrapolants is
$\sim5\cdot10^{-6}$, i.e.\ $\pm0.0006$ in $c_3$, and the $\pm0.01$ absorbs the unverified
assumption that the leading window correction is $\kappa_c^{-2}$ rather than, say,
$\kappa_c^{-2}\log\kappa_c$.  Good practice.

\paragraph{Two things I would ask for.}  (i) The $1/\kappa_c^2$ law for $C(\kappa_c)$ is asserted
and fitted but not derived; since the same document shows (\S4, \texttt{verdictgap}) that the two
constituents $2\Phi_2^{\mathrm{bil}}$ and $\Phi_3$ separately converge only like $1/\kappa_c$, the
$\kappa_c^{-2}$ of the sum is exactly the cancellation whose absence would change the answer.  A
one-line asymptotic model of the cancelling tails would turn $\pm0.01$ into $\pm0.001$.
(ii) The double-precision claim ``the box symbol is accurate to $10^{-14}$, so
$\varepsilon\ge10^{-10}$ is safe'' should be checked against
\texttt{rigor/certified\_numerics.tex}, which finds double precision marginal near the window edge.

\subsection{Does the finite-$\zeta$ table come out?  Yes --- and better than claimed}

RR compares against the \emph{raw} window values of \texttt{results\_A2.txt}
($5.679016,\,2.252719\cdot10^{1},\,8.863062\cdot10^{1},\,3.431959\cdot10^{2}$ in units of
$10^{-7}$) rather than against the published Table~3 of \texttt{fidelity\_tables.tex}, which is
tail-corrected.  I checked that the two differ by exactly the model tail
$\zeta^2/(15\kappa_c^2)$, $\kappa_c=18.42$: raw${}+{}$tail ${}=5.815464,\,23.07298,\,90.81379,\,
351.9286$ against Table~3 column~A $5.8155,\,23.073,\,90.814,\,351.93$ (units $10^{-7}$), agreeing
to the last printed digit.  So the comparison is legitimate, and RR's arithmetic
($\Phi_P/\zeta^2$, the slopes, the implied $c_4^{(P)}\approx+0.87$ stable over a factor $8$ in
$\zeta$) reproduces on my machine to the digits printed.

Then I did the test the brief actually asks for, against the \emph{published} table:
\begin{center}\small
\begin{tabular}{rlll}\toprule
$\zeta$ & Table~3 col.~A & $f_2\zeta^2(1+c_3\zeta+c_4\zeta^2)$, $c_3=-0.99$, $c_4=+0.9$ & ratio\\
\midrule
$1/120$ & $5.8155\cdot10^{-7}$ & $5.815487\cdot10^{-7}$ & $0.999998$\\
$1/60$ & $2.3073\cdot10^{-6}$ & $2.307285\cdot10^{-6}$ & $0.999994$\\
$1/30$ & $9.0814\cdot10^{-6}$ & $9.081380\cdot10^{-6}$ & $0.999998$\\
$1/15$ & $3.5193\cdot10^{-5}$ & $3.519973\cdot10^{-5}$ & $1.000191$\\
$2/15$ & $1.3243\cdot10^{-4}$ & $1.326932\cdot10^{-4}$ & $1.001988$\\
\bottomrule\end{tabular}\end{center}
\label{tab:tab3}
For comparison, $c_3=-0.96$, $c_4=0$ (the brief's reading) is off by $1.9\cdot10^{-4}$,
$2.5\cdot10^{-4}$, $-2\cdot10^{-6}$, $-1.9\cdot10^{-3}$, $-1.2\cdot10^{-2}$ --- visibly a worse fit
with a systematic trend --- and $c_3=-1$, $c_4=+1$ is off by $\sim10^{-4}$ at the small-$\zeta$ end.
The published table therefore \emph{selects} $c_3=-0.990$ to about $\pm0.002$ if one grants the
tail model, and the four-parameter-free prediction reproduces five significant digits over a factor
$16$ in $\zeta$.  This is a much stronger statement than the one made in \S3.2 of the paper and
should be in it.

\begin{verdictok}
$c_3=-0.99(1)$, $c_4=+0.9(1)$: reproduced, and independently corroborated by the published
Table~3.  Caveat, which the paper should state: the six-digit agreement above is \emph{not} an
independent confirmation of the tail model, because the same model tail relates the raw window
values to Table~3.  What it does confirm is that the discarded tail is purely quadratic in $\zeta$
to high accuracy: $f_2-f_2^{(P)}=1.966\cdot10^{-4}$ against $1/(15\kappa_c^2)=1.965\cdot10^{-4}$,
while $f_2c_3-f_2^{(P)}c_3^{(P)}=-1.6\cdot10^{-5}$, i.e.\ the tail's own cubic part is $0.2\%$ of
$f_2c_3$.  Note also that \texttt{check\_note5\_numerics.tex} estimates the error of the tail
correction at $12$--$56\%$ (Note~5, \S6), which would be $\pm1.3\%$ on $\Phi(1/120)$; the observed
$2\cdot10^{-6}$ agreement is far better than that pessimistic bound and should not be
over-interpreted until the discrepancy between the two error estimates is understood.
\end{verdictok}

\section{(c) The lower bound with a rate}\label{sec:c}

\subsection{Proposition 4.1 is correct}
I checked the three optimisations of $h(\Lambda)=7.60\Lambda^{-2}+C(\Lambda)\zeta/f_2$.
For $C=Ce^\Lambda$: $h'=0$ gives $e^\Lambda\Lambda^3=15.2f_2/(C\zeta)$, so
$\Lambda=\log\frac1\zeta-3\log\log\frac1\zeta+O(1)$ and the second term at the optimum is
$15.2/\Lambda^3=o(\Lambda^{-2})$, leaving $7.60\log^{-2}(1/\zeta)(1+o(1))$.  For $C=c\Lambda^m$:
$\Lambda_*^{m+2}=15.2f_2/(mc\zeta)$ and $h(\Lambda_*)=\Theta(\Lambda_*^{-2})
=\Theta(\zeta^{2/(m+2)})$; in particular $m=1$ gives $\zeta^{2/3}$, as used in the
\texttt{verdictgap}.  Bounded $C$ gives $C_\infty\zeta/f_2$.  Arithmetic and asymptotics correct.
The constant $7.60=0.513/g_B=0.513/(8f_2)$ checks out.

\subsection{Proposition 4.2 and its \texttt{verdictbreak}: mis-scoped}

\begin{verdictbreak}
This is my one real objection to the mathematics of the document.  Proposition~4.2 concludes from
the boundedness of the \emph{third Taylor coefficient} that ``the guess $C(\Lambda)\sim e^\Lambda$
of \texttt{quadratic\_limit.tex} Problem~7.2 \ldots is \emph{too pessimistic}''.  It is not, for
Problem~7.2 \emph{as stated}, which is the two-sided bound
$\bigl|\Phi_P(\zeta)-\frac{\zeta^2}8g_B^{(P)}\bigr|\le C(\Lambda)\zeta^3$ on a fixed interval
$0\le\zeta\le\zeta_0$.  Such a $C$ must dominate the quartic Taylor coefficient times $\zeta_0$, and
that coefficient \emph{does} grow exponentially in the window cutoff.  Reason: the exact
$\zeta^4$ coefficient of $\Phi_P$ contains the second-order (quantum Fisher) weight of the
\emph{second}-order defect,
\[
 \tfrac14\sum_{\abs\kappa<\kappa_c}\frac{\abs{(\mathcal D_2)_{\kappa\kappa}}^2}{N_{\kappa\kappa}},
 \qquad (\mathcal D_2)_{\kappa\kappa}=\tfrac2{15}\abs\kappa^{-3},\quad
 N_{\kappa\kappa}=2q_\kappa(1-q_\kappa)\approx2e^{-\abs\kappa},
\]
which is $\approx\frac1{225}e^{\kappa_c}\kappa_c^{-6}$ --- this is exactly the divergence N5
Cor.~4.2 diagnoses, and it is a statement about an honest derivative at $\zeta=0$, unaffected by the
saturation that cures it at finite $\zeta$.  Numerically (my \texttt{scipy} evaluation of the
integral with the exact $q_\kappa=1/(1+e^{\abs\kappa})$): $0.0056,\,0.0067,\,0.0233,\,0.0872,
\,0.4206,\,477$ at $\kappa_c=9.21,\,13.82,\,18.42,\,20.72,\,23.03,\,32.2$.  RR's own measured
quartic coefficient at $\kappa_c=18.42$ is $c_4^{(P)}f_2^{(P)}=0.0072$, the same order.  So
$C(\Lambda)$ in the two-sided sense grows at least like $e^\Lambda/(225\Lambda^6)$: P2's guess is
right, and Proposition~4.2's conclusion ``it is neither exponential nor polynomial in $\Lambda$: it
is bounded'' is false for the quantity Problem~7.2 defines.
\end{verdictbreak}

\begin{verdictgap}
The repair is easy and makes the section \emph{better}, not worse.  Proposition~4.1
$\langle1\rangle1$ only ever uses the \emph{one-sided} inequality
\begin{equation}\label{eq:onesided}
 \Phi_P(\zeta)\ \ge\ \tfrac{\zeta^2}8g_B^{(P)}-C^-(\Lambda)\,\zeta^3,\qquad 0\le\zeta\le\zeta_0 ,
\end{equation}
and for the one-sided constant the exponentially large quartic term has the \emph{helpful} sign
($c_4^{(P)}>0$), so it does not force $C^-$ to grow: with $c_3^{(P)}<0<c_4^{(P)}$ the function
$-f_2^{(P)}(c_3^{(P)}+c_4^{(P)}\zeta+\cdots)$ is decreasing in $\zeta$ and its supremum on
$[0,\zeta_0]$ sits at $\zeta=0$, i.e.\ at the third Taylor coefficient --- which is what
Proposition~4.2 measures.  So: (i) restate Problem~7.2 one-sidedly as \eqref{eq:onesided};
(ii) keep Proposition~4.2 as evidence for $C^-(\Lambda)=O(1)$, i.e.\ for the third row of
Proposition~4.1; (iii) delete or completely rewrite the \texttt{verdictbreak} against P2.  Two
further gaps then remain, both correctly identified by RR: the measured quantity is a Taylor
coefficient, whereas \eqref{eq:onesided} needs the inequality on an interval uniformly in $\Lambda$;
and the ingredients (i)--(iii) of RR's own \texttt{verdictgap} (exponential decay of
$\mathcal D_1$, a two-variable bound on $\mathcal D_2$, and the cancellation of the two
$1/\kappa_c$ tails) are the analytic content.

A cheap and decisive numerical test the author should run: measure $c_4^{(P)}$ at
$\kappa_c=9.21,\dots,23.03$ with the machinery already in \texttt{rr\_c3c.py}.  My model predicts
it grows by a factor $\approx e^{\Delta\kappa_c}(\kappa_c/\kappa_c')^6$, i.e.\ by $\sim18$ between
$\kappa_c=18.42$ and $23.03$, while $c_3^{(P)}$ stays put.  If that is what happens, the diagnosis
above is confirmed and the one-sided reformulation is forced.
\end{verdictgap}

\section{(d) Is the refutation of the $\log^{-2}$ remainder logically complete?}\label{sec:logic}

Let me separate the four statements that \S5 of \texttt{rate\_of\_remainder.tex} runs together.

\begin{enumerate}[label=(R\arabic*),leftmargin=3em]
\item\label{R1} \emph{$\Phi(\zeta)-f_2\zeta^2\le f_2^2\zeta^4/(1-2f_2\zeta^2)$ for all
 $\zeta<(2f_2)^{-1/2}$, hence
 $\limsup_{\zeta\downarrow0}\log^2(1/\zeta)\,\zeta^{-2}\bigl(\Phi(\zeta)-f_2\zeta^2\bigr)\le0$.}
 \textbf{Proved, conditional on $f_2^\star=f_2$} --- and unconditionally proved if
 Proposition~\ref{prop:repair} is adopted.  This is a complete refutation of N5 Cor.~4.2's
 remainder \emph{as a net statement about $\Phi$}: the net deviation of $\Phi$ from $f_2\zeta^2$
 cannot be a positive quantity of order $\zeta^2/\log^2(1/\zeta)$.
\item\label{R2} \emph{The unconditional variant.}  Cor.~2.7 $\langle1\rangle2$ excludes N5's term
 only on $\zeta\ge6\cdot10^{-17}$ (resp.\ $2\cdot10^{-23}$).  This is \emph{not} a refutation of an
 asymptotic ($\zeta\downarrow0$) claim, only of the numerically relevant regime.  The document is
 honest about this in the corollary but the abstract's flat ``\emph{refuted}'' is only justified
 once $f_2^\star=f_2$ is in hand.  One more reason to adopt Proposition~\ref{prop:repair}.
\item\label{R3} \emph{$c_3\le0$, and $c_3=-0.99(1)$.}  The inequality is a consequence of
 \ref{R1} \emph{plus} the assumption that the limit $c_3=\lim(\Phi/(f_2\zeta^2)-1)/\zeta$ exists.
 That the limit exists, and its value, are numerical (\S\ref{sec:c3num}); the paper says so.
\item\label{R4} \emph{The mechanism} (\S5.1(i)--(ii): the UV saturation is a re-count of the share
 of $f_2$ already carried by the modes above $\kappa_*$, and the two-mode rotation example).
 \textbf{Heuristic.}  The numerical coincidence $\frac{\zeta^2}{15\kappa_*^2}$ versus
 $\frac{\zeta^2}8(g_B-g_B^{(\kappa_*)})=\frac{\zeta^2}{15.6\kappa_*^2}$ is striking and, I agree,
 is the explanation; but ``these agree to $4\%$, so the remainder is the difference of two nearly
 equal quantities'' is a plausibility argument, not a proof, and $4\%$ of
 $\frac{\zeta^2}{15\kappa_*^2}$ is still $\gg\zeta^4$.  \ref{R1} is what does the work.
\end{enumerate}

\begin{verdictok}
The refutation is logically complete at the level of \ref{R1}, which is the statement that matters:
an upper bound with no term between $\zeta^2$ and $\zeta^4$ is incompatible with a positive
$\zeta^2/\log^2(1/\zeta)$ remainder.  The proposed replacement text for N5 Cor.~4.2 (\S5.2) is
accurate in substance.  Three edits I would insist on: (1) mark which of $f_2$, $c_3$, $c_4$,
``integer powers'', ``$c_3$ independent of $r$'' are theorems and which are numerics --- as written,
the italic block reads as a theorem and only its last two lines are proved;
(2) drop ``$c_3$ is independent of $r$'' or justify it (it follows from the $r$-fold direct-sum
structure of PA Thm.~8.1 $\langle1\rangle1$, which multiplies $\Phi$ and $f_2$ by the same $r$ ---
one line, but it should be there); (3) the retained sentence about the divergent formal fourth-order
coefficient should say \emph{how} it diverges, since \S\ref{sec:c} shows that its rate
($e^{\kappa_c}\kappa_c^{-6}$ in the compressed problem) is exactly what makes the two-sided
Problem~7.2 constant exponential.
\end{verdictok}

\section{Citations, typos, internal inconsistencies}\label{sec:minor}

\paragraph{Cross-references to PA are systematically shifted by one section.}  Against the current
\texttt{uhlmann\_upper\_bound.tex} (I read the numbers off \texttt{uhlmann\_upper\_bound.aux}):
\begin{center}\small
\begin{tabular}{llll}\toprule
cited as & should be & cited as & should be\\\midrule
PA Lemma 2.3 (kernel) & Lemma 2.4 & PA Lemma 4.2 (integral rep.) & Lemma 5.2\\
PA Prop.\ 2.4 (HS bound) & Prop.\ 2.5 & PA Def.\ 4.4 (admissible $h'$) & Def.\ 5.3\\
PA Prop.\ 2.7 (Uhlmann) & Prop.\ 3.4 & PA Prop.\ 4.5 (expansion) & Prop.\ 5.4\\
PA Thm.\ 3.1 (overlap) & Thm.\ 4.1 & PA Lemma 6.3 ($H_{\cM'}$) & Lemma 7.5 / Prop.\ 7.6\\
PA Lemma 4.1 (generator) & Lemma 5.1 & PA Thm.\ 7.1, 7.2 & Thm.\ 8.1, 8.2\\
\bottomrule\end{tabular}\end{center}
(``PA Prop.\ 2.7'' currently names \emph{Corollary 2.7, The implementer}, and ``PA Thm.\ 7.2''
names \emph{Value of the quasi-free supremum} --- both real but different results, so the reader is
actively misled.)  Everything cited is present in PA under the corrected number, and in every case
I verified that the \emph{content} used is the content of the intended result.

\paragraph{Other items.}
\begin{itemize}\itemsep1pt
\item \textbf{Abstract vs.\ body, twice.}  The abstract says ``$c_3=-0.96(3)$''; eq.~(16) and \S5.2
 say $-0.99(1)$.  The abstract says $C(\Lambda)$ ``grows \emph{polynomially}, not exponentially, so
 the optimised lower rate is a \emph{power}''; Prop.~4.2 and the summary table say it is
 \emph{bounded} and the rate is \emph{linear}.  Both must be reconciled (and, per \S\ref{sec:c},
 in the two-sided reading neither is right).
\item \textbf{Sign typo.}  In \S3.2, ``$c_4^{(P)}:=(c_3^{(P)}-\text{slope})/\zeta$'' should be
 $(\text{slope}-c_3^{(P)})/\zeta$; with the printed definition the table's $+0.874$ would be
 $-0.874$.  The table is right, the definition is wrong.
\item \textbf{``for all $s$''} after (F2) should be ``for all $s\ge0$'' (\S\ref{sec:a1}).
\item \textbf{Theorem 2.3 $\langle1\rangle1$} cites a numerics output file
 (\texttt{p1\_bogoliubov.out}~(1)) for a mathematical identity.  Replace by the two-line proof
 (Theorem~\ref{thm:mine} $\langle1\rangle1$--$\langle1\rangle2$).
\item \textbf{N5 numbering.}  The brief calls the quasi-free fidelity formula ``Note 5 Thm.~3.2'';
 the paper cites it as ``N5 Thm.~2.1'' via \texttt{check\_note5\_theory.tex}.  Fix one of the two.
\item \textbf{Prop.~2.6 $\langle1\rangle2$} should record that each finite-$M$ minimiser in the
 basis $\sin^2(\pi u)\{\cos,\sin\}(m\pi u)$ really is an \emph{admissible} field in the sense of
 Definition~2.1 --- it is (the second-order vanishing at the arc ends makes $\chi$ tend to a
 nonzero constant, not grow, at $x\to+\infty$, so $\kappa^{(1)}\in L^2$) --- because that is the
 whole content of ``each finite $M$ is an admissible smooth field''.
\end{itemize}

\section{Recommendation and status board}

\textbf{Recommendation: accept after revision.}  The document is correct where it claims to be
proved, is scrupulous about labelling its numerics, and its central observation --- that a
one-parameter \emph{group} turns PA's Step~4 from an asymptotic expansion into an exact inequality
with no error term --- is right, is important, and removes an entire class of remainder terms from
the problem at one stroke.  Required revisions, in order of importance:

\begin{enumerate}[leftmargin=2.2em]
\item Adopt Proposition~\ref{prop:repair} (or refute it): the group $e^{s(D+ih')}$, $h'\in\cA'$, is
 admissible and makes $f_2^\star=f_2$ a theorem, removing the conditional clause from the main
 result and from the refutation of N5 Cor.~4.2.
\item Rewrite \S4: Problem~7.2's constant is exponential in the two-sided reading
 (\S\ref{sec:c}); state and use the one-sided version, and delete the \texttt{verdictbreak} against
 \texttt{quadratic\_limit.tex}.  Measure $c_4^{(P)}$ versus $\kappa_c$.
\item Fix the PA cross-references (\S\ref{sec:minor}) and the two abstract/body inconsistencies.
\item Replace the numerics-file citation in Thm.~2.3 $\langle1\rangle1$ by the two-line proof; add
 the $s\ge0$ qualifier; separate theorem from conjecture in the proposed N5 Cor.~4.2 replacement.
\end{enumerate}

\begin{center}\small
\begin{tabular}{p{0.46\textwidth}l}\toprule
statement & status after this report\\\midrule
$\norm{V_s}_2^2\le s^2\norm{\Pm K\Pp}_2^2$ for any one-parameter unitary group with
 $[\Pm,K]\in\Stwo$ & \textbf{proved} (stronger than Thm.~2.3)\\
$\Phi(\zeta)\le-\frac12\log(1-2f_2^\chi\zeta^2)$, $\chi$ an admissible field &
 \textbf{proved} (Thm.~2.5)\\
$f_2^\star\ge f_2$ & \textbf{proved}\\
$f_2^\star=f_2$, hence $\Phi(\zeta)\le-\frac12\log(1-2f_2\zeta^2)=f_2\zeta^2(1+f_2\zeta^2+\dots)$ &
 numerical in the paper; \textbf{provable} by Prop.~\ref{prop:repair}\\
$\limsup\log^2(1/\zeta)\zeta^{-2}(\Phi-f_2\zeta^2)\le0$: no positive $\zeta^2/\log^2$ remainder &
 same conditionality; \textbf{refutation of N5 Cor.~4.2 stands}\\
$c_3\le0$ & follows, given that $c_3$ exists\\
$\det(\one-V_s^*V_s)$ even in $s$ & \textbf{proved} (and not needed)\\
third-order formula (3.3) for $-\log\Fid$ & \textbf{proved} and numerically verified\\
$c_3=-0.99(1)$, $c_4=+0.9(1)$ & numerical; corroborated to $2\cdot10^{-6}$ by Note~5 Tab.~3 col.~A\\
$C^-(\Lambda)$ bounded $\Rightarrow$ $\Phi\ge f_2\zeta^2(1-0.99\zeta+O(\zeta^2))$ &
 conditional; the Taylor coefficient is measured, the uniform bound is not\\
$C(\Lambda)$ (two-sided, Problem 7.2) is bounded & \textbf{false}: it grows like
 $e^\Lambda\Lambda^{-6}$\\
\bottomrule\end{tabular}\end{center}
\end{document}
